% Chapter 6: Newton's Three Rules (complete sample chapter)
\chapter{Newton's Three Rules}

\step{A Counterintuitive Opening}

\begin{mainline}
On a Winter Olympics curling sheet, an athlete pushes a granite stone of about \(19\) kilograms forward and releases it. From that instant nobody touches it. No engine, no rope, yet it slides more than thirty meters until it hits another stone or the ice slowly brings it to a stop.

Change the scene: put the same stone on grass and push with about the same effort. How far does it slide? Less than half a meter.

The same push-and-release gives results differing by dozens of times. The grass is not the part that should keep you awake; the ice is. After release, \emph{who} keeps the stone going? Every part of our intuition says that moving things need a continuous push. When someone stops pushing a cart, it stops. Isn't that obvious? Yet the ice contradicts intuition: less friction means less need for anyone to push, and the object goes straighter and farther on its own. What would a stone do on perfectly smooth, infinitely long ice?

Push the scene farther. Outside a spacecraft, an astronaut loses their grip on a wrench while tightening a bolt, and it drifts away slowly. No rocket, no wings, not even air nearby, yet it travels straight, neither speeding up, slowing down, nor turning until it hits something. Who maintains its motion? Nobody. That nobody is this chapter's starting point.

Friction and air surround everyday life, burying a profound rule beneath our experience that pushing a cart requires continued effort. This chapter scrapes away that thick experience to uncover three almost brutally clean rules: Newton's laws of motion. They govern curling stones, wrenches, every forward lurch on a bus, and rockets lifting off.
\end{mainline}

\step{Make a Guess First}

Before reading on, treat this page as a betting table and stake your intuition. Do not peek at the answers. Wrong guesses are more valuable than right ones, because intuition is what we aim to repair.

Question one. Imagine perfectly smooth, endless ice with absolutely no friction. You gently push a curling stone and release it. What follows? (a) It slows gradually and stops; (b) it keeps sliding uniformly in a straight line forever; (c) it gets faster and faster. Write your choice and reason.

Question two. Lay a sturdy paper strip on a table with a plastic cup half full of water on top and one end hanging over the edge. Pull the strip horizontally out as hard and fast as you can. Does the cup (a) fly away with it, (b) topple in place, or (c) hardly move? What if you pull slowly instead?

Question three. Standing in a small boat on a lake, you push water backward hard with an oar. The water goes backward and the boat forward. But however hard you push the water, it pushes you equally hard. If those forces are equal and opposite, should they not cancel? How can the boat advance?

Write your three answers first. At the chapter's end we will check each. All correct is good; wrong guesses are even better, because you will soon see exactly where your intuition turned the wrong way.

\step{Hands-on Experiment}

\begin{experiment}[Pulling Paper: Why the Cup Is Reluctant to Move]
\textbf{Materials}: a sturdy paper strip, cut from newspaper or kitchen paper, about five centimeters wide and thirty long; a plastic cup; water; and a table as smooth as possible.

\textbf{Step one}: half-fill the cup and stand it on the strip, leaving about ten centimeters hanging over the edge. Pinch the end and pull horizontally \emph{slowly and steadily}. The cup nearly follows the paper. This is question two's slow-pull case: friction has ample time to drag the cup along.

\textbf{Step two}: return the cup to its starting position. Grip the strip and jerk it out horizontally as fast as possible. Keep the movement quick and level, without lifting. The paper whips out and the cup only trembles, barely moving.

\textbf{Step three}: fill the cup completely, making it heavier, and jerk again. The cup is steadier. Try a lighter object, such as an empty paper cup: it is carried along more easily.

\textbf{Record}: why do slow pulling and sudden jerking produce such different behavior? What about heavy and light cups? Record observations without explaining them yet. We will unpack the whole secret gradually: the key difference is not force magnitude, but \emph{how long} it acts. A heavier cup is more reluctant, a trait we will formally name inertia.
\end{experiment}

Another free observation: next time you ride a bus or subway, stand securely and hold on. Note your body's leaning direction when it \emph{starts} and \emph{brakes}, recording three examples. Braking makes you lean forward; starting makes you lean back. Most people say a force throws them forward when braking. Remember that feeling: later we will show that this forward force does not exist.

\step{The Old Model Runs into Trouble}

First give intuition its due. More than two thousand years ago, Aristotle summarized everyday experience in an apparently unassailable account: for an object to move, something must push it; remove the push and motion stops. This old model successfully explained pushing carts, rowing, and turning millstones for so long that almost nobody questioned it for two millennia. It is exactly the intuition we have accumulated since pushing toy cars in infancy.

Consider why it lived so long: we happen to inhabit a corner of the universe filled with friction and air resistance. Here, moving objects really do stop by themselves, making removal of force followed by stopping a dependable experience. The old model is not stupidity; it is a faithful but unexamined observation report. Its mistake is confusing the usual conditions in our corner with universal law. To expose it, find settings with sufficiently little friction: ice, air cushions, and space supply them.

But difficulties arrive one after another.

Difficulty one: the curling stone and space wrench. Release removes the push but not the motion. The model says motion should stop without force, yet the stone slides another thirty meters. It cannot answer who pushes over those thirty meters.

Difficulty two: Galileo's inclined-plane reasoning. A ball rolling down a slope speeds up; rolling up the opposite slope it slows, but nearly regains its original height. Make the opposite slope gentler and it must travel farther before stopping. Smoother surfaces bring it closer to the original height and carry it farther. Now take the famous limit: lower the opposite slope until perfectly horizontal and imagine it frictionless. Can the ball ever regain its original height? Never, so it must keep rolling horizontally, neither speeding up nor slowing. The argument's sharpness is that you need not build a truly frictionless world: progressively reduce friction and watch where the trend points. It is unmistakable: \emph{less friction means less need for anything to maintain motion}. The old model demands force to sustain motion, while experiments point to motion being most stable precisely as force vanishes.

Difficulty three: the directions disagree. Throw a ball straight up. After release its force, gravity, points downward, yet it keeps traveling upward for a while before turning. If force maintained motion, how could downward force permit continued upward motion?

All three share one diagnosis: the old model combines two separate things, maintaining motion and changing motion. The stone keeps moving not because something sustains it, but because nothing changes it. Upward flight needs no upward force; only velocity's \emph{change} concerns gravity. Separating these requires new roles and concepts obvious now but earth-shattering when introduced.

\step{Inventing New Concepts}

\begin{mainline}
\textbf{First concept: inertia (Newton's first law).} Galileo's slopes show that a ball with no horizontal push or pull rolls uniformly in a straight line forever. Newton made this his first rule: \emph{when external forces on an object cancel, giving zero net force, it remains at rest or in uniform straight-line motion}. This reluctance to change its state of motion is inertia. Inertia is \emph{not} a force. Nobody is pushing from behind; it is simply the fact that motion stays unchanged without external intervention. The first law also does something major: it places rest and uniform straight-line motion in one class. Both are states requiring no explanation; only \emph{change} needs explaining.

\textbf{Second concept: force is interaction, not a stockpile.} The old model imagines force stored inside objects like sweets in a pocket: the pusher gives force to a cart, it travels carrying force, and stops when force runs out. This picture is entirely wrong. The new definition is \emph{that force always represents interaction between two objects}, one pushing and another being pushed, both essential. What is it like? A handshake, which requires two hands: you cannot shake hands with yourself or pocket a handshake for later. What is it unlike? Water in a jug: force cannot be stored inside an object or used up. How much force an object contains is therefore an ill-formed question. An object can possess velocity, energy, and momentum, the subjects of Chapters 9 and 10, but force is only the interaction it is \emph{currently} having with its surroundings. After release, no forward force acts on the stone. This does not impede its progress; it is precisely why it moves straight and steadily.

\textbf{Third concept: mass is resistance to velocity change.} An empty shopping cart shoots forward at a gentle push; a cart full of rice starts sluggishly under the same effort. The difference is greater resistance to velocity change. We call that degree of resistance \emph{inertial mass}, or mass. Greater mass means less velocity change from the same force. Be especially careful: mass is not weight. Weight is Earth's force on the object, detailed in Chapter 8; mass is its intrinsic resistance to acceleration. Move astronaut and cart to a weightless space station and weight disappears, but the rice-filled cart remains sluggish to push. Not an ounce of mass has vanished.

\textbf{Fourth rule: net force determines how motion changes (Newton's second law).} Since motion itself needs no explanation, change does. Newton says the rate of velocity change, acceleration, is determined by \emph{net force} and points in its direction. For equal net force, greater mass means smaller acceleration. More fundamentally, \emph{net force determines the rate of change of momentum}. Remember that statement now; Chapter 10 will make it central. The formula specifies net force: a book on a table experiences gravity and support, which cancel to zero, so it stays still. It rests not because no forces act, but because their total is zero.

\textbf{Fifth rule: forces always come in pairs, acting on different objects (Newton's third law).} Push a wall and it pushes you; an oar pushes water backward and water pushes the oar forward; Earth pulls an apple down and the apple pulls Earth up. The pair is equal and opposite, but \emph{acts on different objects}, so it can never be canceled against itself. Cancellation occurs only among forces on the same object.

Accept paired forces and daily life becomes an exhibition of the third law. Walking: your foot pushes the ground backward and the ground pushes you forward. Fail to get that purchase on smooth ice and you immediately understand the source of forward force. Swimming: arms push water backward and water carries you forward. A hand hurts on hitting a wall because a harder push receives a harder push back; it is not merely hardness, but meticulous obedience to the third law. Gun recoil strikes a shoulder backward: propellant gas pushes the bullet forward and the bullet pushes the gun backward. These diverse examples have one skeleton: to propel yourself forward, push something else backward.
\end{mainline}

The new concepts now dismantle a famous objection, the opening third question and a stumbling block for countless beginners: \textbf{the counterintuitive horse-and-cart paradox}. A thoughtful horse tells Newton: ``Your third law says my pull on the cart equals its opposite pull on me. Their sum is always zero, so how can the cart move? I might as well stop pulling.'' Its mistake is adding forces on different objects. To judge the \emph{cart}, consider only its forces: the horse pulls forward and ground resistance drags backward. A greater forward force accelerates it forward. To judge the \emph{horse}, consider only its forces: the cart pulls back, but the hooves push backward on the ground, which pushes forward through static friction. If that push exceeds the cart's pull, the horse accelerates forward. Those equal, opposite partners never act on the same recipient, so why speak of cancellation?

\begin{center}
\begin{tikzpicture}[scale=1.0]
  % Ground
  \draw[thick] (-4.8,0) -- (4.8,0);
  \fill[pattern=north east lines] (-4.8,-0.28) rectangle (4.8,0);
  % Cart (left), horse (right), motion to the right
  \draw[thick,fill=gray!15] (-3.0,0) rectangle (-1.4,1.0);
  \node at (-2.2,0.5) {Cart};
  \draw[thick,fill=orange!15] (1.4,0) rectangle (3.0,1.0);
  \node at (2.2,0.5) {Horse};
  \draw[thick] (-1.4,0.5) -- (1.4,0.5) node[midway,above=1pt] {Rope};
  % Forces on cart
  \draw[-{Stealth[length=3mm]},very thick,red!70!black] (-2.2,1.25) -- (-0.5,1.25) node[right=2pt] {Horse pulls cart (forward)};
  \draw[-{Stealth[length=3mm]},very thick,blue!70!black] (-2.2,0.5) -- (-3.9,0.5) node[left=2pt] {Ground resistance};
  % Forces on horse
  \draw[-{Stealth[length=3mm]},very thick,red!70!black] (2.2,1.55) -- (0.5,1.55) node[left=2pt] {Cart pulls horse (backward)};
  \draw[-{Stealth[length=3mm]},very thick,blue!70!black] (2.2,0.5) -- (3.9,0.5) node[right=2pt] {Ground pushes horse};
\end{tikzpicture}
\end{center}

This diagram also introduces the chapter's final tool: \textbf{free-body diagrams and system boundaries}. There are two steps: draw an imaginary dashed boundary around the object of study---cart, horse, or both together---then draw \emph{only external forces acting on it}, nothing else. Change the boundary and the problem changes appearance. For horse plus cart as one system, their mutual pulls become internal forces and need not be considered. The system advances through ground friction pushing the hooves forward minus ground resistance. A correct free-body diagram completes half a mechanics problem. An incorrect one, such as inserting forces acting on other objects, cannot be rescued by skillful calculation.

\step{The Model in One Sentence}

\begin{oneline}
Force determines how fast your velocity changes, not how fast you move now; forces always occur in pairs---however hard you push the world, it pushes you equally hard, with the two forces acting on two different objects.
\end{oneline}

This contains all three laws. The first half gives laws one and two: unchanging motion is the default, and net force creates change. The second half gives law three. Remember it and this chapter will not have been wasted.

\step{The Mathematical Translator}

\begin{mainline}
Translate the previous section into one line:
\[
  F_{\text{net}} = ma
\]
Reading it matters more than writing it. On the left, \(F_{\text{net}}\) is the total of all external pushes and pulls on this object, asking how its surroundings collectively seek to change it now. On the right, \(a\) is acceleration, asking how rapidly velocity actually changes. The \(m\) belongs in the denominator, more clearly seen in \(a=F_{\text{net}}/m\): greater mass buys less acceleration for the same net force. That is the mathematical face of mass as resistance to acceleration.

Three routine checks come first. Set \(F_{\text{net}}=0\), obtaining \(a=0\): constant velocity. The first law appears automatically, showing the three rules belong to one model, not three unrelated slogans. Double \(m\), and \(a\) halves: the loaded shopping cart starts half as quickly, matching experience. Check direction: \(a\) always follows \(F_{\text{net}}\), but may oppose \emph{velocity}. Braking gives backward net force while the car still travels forward and slows. The old model's most frightening scene becomes ordinary arithmetic in the new one.

Practice numerical intuition. An empty cart has mass about \(15\) kilograms. Push horizontally with \(30\) newtons, roughly the effort of supporting three one-liter cola bottles, and neglect friction: \(a=30/15=2\) meters per second squared. Starting from rest, one second gives \(2\) meters per second, two seconds \(4\). No wonder an empty cart jumps at a touch. Fill it with rice to a total \(60\) kilograms and the same \(30\) newtons buys only \(0.5\) meters per second squared: four times the resistance, one-quarter the result. That is \(a=F_{\text{net}}/m\) in the supermarket aisle.
\end{mainline}

Now estimate with real data to see how the formula governs survival. A car hits a wall at about \(50\) kilometers per hour, roughly \(14\) meters per second. An unbelted passenger strikes the dashboard almost instantly, slowing in about \(0.05\) seconds from \(14\) meters per second to zero. Their acceleration magnitude is about
\[
  a = \frac{14\ \text{m/s}}{0.05\ \text{s}} = 280\ \text{m/s}^2,
\]
nearly thirty times gravitational acceleration. A \(60\)-kilogram person experiences \(F=ma\approx 60\times 280 \approx 17000\) newtons, like a car's weight pressing on them. With a seat belt and airbag, the belt stretches and the bag vents, extending deceleration to about \(0.5\) seconds. Acceleration falls to about \(28\) meters per second and force to about \(1700\) newtons, roughly three adults' weight on the chest: painful but usually survivable. Nothing about the velocity change has shrunk; \(F_{\text{net}}=ma\) simply spreads the same account over longer time. Engineers design seat belts, airbags, bicycle helmets, and parcel foam using this line. The idea that a given velocity change over more time requires less force will become Chapter 10's impulse--momentum theorem. Remember its shape for now.

\begin{mainline}
See where the free-body diagram enters the formula. For a book on a table, the diagram below shows Earth pulling down through gravity and the table pushing up through its normal force. Both act on the \emph{book}, equal and opposite, giving zero net force, so \(a=0\) and the book rests. Here is the chapter's easiest trap: these are \emph{not} an action--reaction pair! They act on the same object, whereas third-law partners must act on two objects. Gravity's partner is the book pulling Earth upward; the normal force's partner is the book pushing the table down. Equilibrium, zero net force on one object, and the third law, paired forces between objects, are distinct things that often happen together.
\end{mainline}

\begin{center}
\begin{tikzpicture}[scale=1.0]
  % Tabletop
  \draw[thick] (-2.6,0) -- (2.6,0);
  \fill[pattern=horizontal lines] (-2.6,-0.28) rectangle (2.6,0);
  % Book
  \draw[thick,fill=gray!15] (-0.9,0) rectangle (0.9,0.7);
  \node at (0,0.35) {Book};
  % Both forces drawn on the book
  \draw[-{Stealth[length=3mm]},very thick,blue!70!black] (0.25,0.35) -- (0.25,1.85);
  \node[right=2pt] at (0.25,1.85) {Normal force (table pushes book)};
  \draw[-{Stealth[length=3mm]},very thick,red!70!black] (-0.25,0.35) -- (-0.25,-1.15);
  \node[right=2pt] at (-0.25,-1.15) {Gravity (Earth pulls book)};
\end{tikzpicture}
\end{center}

\begin{mathbridge}
Force has direction, so the full second law is a vector equation:
\[
  \vec{F}_{\text{net}} = m\vec{a},
  \qquad\text{where}\qquad
  \vec{F}_{\text{net}} = \vec{F}_1 + \vec{F}_2 + \cdots
\]
In practice, resolve along perpendicular directions: horizontally \(F_{\text{net},x}=ma_x\), vertically \(F_{\text{net},y}=ma_y\), calculating independently. The tabletop book illustrates \(F_{\text{net},y}=0\): support and gravity cancel, giving zero vertical acceleration.

Newton himself actually wrote not \(F=ma\), but the more fundamental form
\[
  \vec{F}_{\text{net}} = \frac{d\vec{p}}{dt},
  \qquad \vec{p}=m\vec{v}\ \text{(momentum)},
\]
net force equals the rate of momentum change. With constant mass, \(\frac{d\vec{p}}{dt}=m\frac{d\vec{v}}{dt}=m\vec{a}\), reducing to the familiar form. It can also handle changing-mass objects, such as a rocket growing lighter as it exhausts gas. Chapter 10 will fully reveal this formula's power.

This also answers a question you may have noticed: the book's opening three questions---what state is the system in, what law changes that state, and how measurement tells us---receive their first complete answers in Newtonian mechanics. Position and velocity give state; \(F_{\text{net}}=dp/dt\) gives its change; measurement uses rulers and clocks to record position over time, while accelerometers and bathroom scales actually measure force. State, law, and measurement will recur throughout the book, with each later part replacing their components.

Finally, use graphs: for one object, \(a\) is proportional to \(F_{\text{net}}\), a straight line through the origin with slope \(1/m\). Greater mass makes the line flatter: the same force buys less acceleration.
\end{mathbridge}

\begin{center}
\begin{tikzpicture}[scale=1.0]
  \draw[-{Stealth[length=3mm]}] (0,0) -- (5.0,0) node[below left] {Net force \(F_{\text{net}}\)};
  \draw[-{Stealth[length=3mm]}] (0,0) -- (0,3.0) node[below left] {Acceleration \(a\)};
  \draw[very thick,blue!70!black] (0,0) -- (4.4,2.2) node[right=2pt] {Small \(m\) (empty cart)};
  \draw[very thick,red!70!black] (0,0) -- (4.4,0.9) node[right=2pt] {Large \(m\) (full of rice)};
  \draw[dashed] (3,0) node[below] {\(F_0\)} -- (3,1.5) -- (0,1.5) node[left] {\(a_0\)};
\end{tikzpicture}
\end{center}

\begin{challenge}
Rockets are the second law's grandest application, shattering the misconception that they rise by pushing against ground or air. They burn carried fuel into gas and expel it rapidly backward: rocket pushes gas back and gas pushes rocket forward. The third law works in vacuum; nothing behind the craft needs to provide support. The complication is continuously decreasing mass: burn half the fuel and the rocket becomes half as heavy, so the same thrust buys ever more acceleration. Apply \(F=\frac{dp}{dt}\) to both rocket and exhausted gas to derive Tsiolkovsky's equation: the eventual velocity increment depends on exhaust speed and the logarithm of initial mass divided by final mass. Logarithms are stingy: a little more speed demands fuel in multiples. That is why launch vehicles have stages, discarding each after burnout. Chapter 10's challenge track will derive the equation; remember only that a rocket in airless space works by throwing part of itself backward.
\end{challenge}

\step{The Model's Limits}

\begin{mainline}
Newton's laws are almost domineeringly useful, but their territory has definite borders. Applying them beyond those borders causes errors. First the conditions, then common misuses.

Condition one: \emph{inertial reference frames}. Newton's laws apply directly only in frames that neither accelerate nor rotate. The ground approximately qualifies, as does a carriage moving uniformly in a straight line. A braking carriage does not: inside it, you see someone surge forward without a horizontal pushing force, impossible to explain directly with Newton's laws in that frame. Introducing inertial forces repairs the bookkeeping, but these are inventions for continuing to use the equations in noninertial frames. The ledger must explicitly note that no force-exerting object can be found for them.

Condition two: \emph{speeds far below light speed}. Near light speed, \(a=F/m\) systematically departs from experiments. Under constant thrust, speed increases ever more slowly and never passes light speed. This is not a minor numerical error: relations between velocity, time, and mass need rewriting in Chapters 38--40. Interestingly, \(F=dp/dt\) outlives \(F=ma\): relativity retains precisely the core that net force changes momentum.

Condition three: \emph{macroscopic objects}. At atomic scales, electrons and atoms no longer follow definite Newtonian trajectories; quantum mechanics, beginning in Chapter 43, is needed. This is not inadequate instruments failing to measure accurately, but the microscopic failure of Newton's premise of simultaneously definite position and velocity.
\end{mainline}

Four misuses specific to this chapter, each matching an earlier trap:

\begin{itemize}[leftmargin=1.8em]
  \item \emph{Objects need force to maintain motion.} Wrong. Net force changes momentum rather than maintaining velocity; uniform straight-line motion needs no force. This fossil of two-thousand-year-old intuition is what the entire chapter dismantles.
  \item \emph{Treating mass as weight.} Mass resists acceleration; weight is Earth's pull. Weight vanishes in space while mass remains, so pushing a large box still takes effort.
  \item \emph{Putting action and reaction on the same object and announcing cancellation.} The horse-and-cart paradox again. Only forces acting on the same object can be added and canceled.
  \item \emph{Circular motion with changing velocity direction has zero net force.} Wrong. Velocity includes direction, and directional change is change. A turning car or rotating swing has nonzero net force. Next chapter will calculate it.
\end{itemize}

\step{Explaining the Opening Phenomenon}

Return to the opening mysteries and reconcile each with the new model.

\textbf{The curling stone.} After release, its only horizontal force is tiny ice friction. A small net force in \(F_{\text{net}}=ma\), or more fundamentally a small rate of momentum change, reduces speed only slowly, allowing over thirty meters of sliding. On grass, much greater friction rapidly drains the same velocity deposit and stops it within half a meter. On perfectly smooth ice the net force is zero, the first law's domain, and the stone slides forever in a straight line at its release velocity. Your opening choice (b) was right: \emph{continuing is the default; stopping needs a reason}. The space wrench is a purer version: with no friction, it drifts straight forever.

\textbf{Pulling the paper.} Pull slowly and friction between paper and cup acts for several seconds, gradually increasing cup velocity so it follows. Jerk quickly and friction is not much greater, but acts for only a few tenths of a second; the paper escapes before much velocity accumulates. The same push acting for less time creates less velocity change: the domestic version of the seat-belt estimate. A full cup has greater mass, so the same frictional push changes velocity less and it remains steadier. You have seen inertia sitting on your dining table.

\textbf{Bus braking.} The carriage slows and friction slows your feet with it, but your upper body receives no backward force. Faithfully obeying the first law, it continues forward at its original velocity and leans forward relative to the carriage. No force throws you forward at any point: in the ground frame, your upper body merely was not held back and continues uniformly. The feeling of being thrown is an illusion of the accelerating carriage frame. Similarly, at startup your feet move forward with the carriage while your upper body lags, making you lean backward. A seat belt supplies the missing backward pull on your upper body so that it decelerates with the carriage instead of flying out.

Three mysteries, one key: separate motion from change in motion. Net force governs only the latter.

\step{Thought Experiments and Challenges}

\begin{thought}[Infinite Ice and a Stone at the Universe's Edge]
Push the imagination to an extreme. On perfectly smooth, infinite ice, give a stone a gentle push, turn off every light in the universe, and forget it. Where is it ten million years later? The first law says still traveling at its original velocity along its original line, at a distance approximately velocity times ten million years. More extreme: a spaceship reaches deep space far from every galaxy, where gravity is negligible, then switches off its engines. Will it slowly stop? No. It will drift uniformly in a straight line at its shutdown velocity forever, after a hundred million or ten billion years alike. Nothing consumes its motion because motion needs no maintenance. This experiment highlights how unusual terrestrial life is: stopping feels natural because our corner is full of friction and drag. The first law is no redundant idealized assumption. It is the world's default, buried under layers of Earth's resistance.
\end{thought}

\begin{thought}[Do Newton's Laws Fail in a Falling Elevator?]
Another thought experiment connects this chapter to distant theories. You stand on a bathroom scale in an elevator when the cable suddenly breaks, and both you and it fall freely under gravity. The scale reads zero; release an apple and it floats before you rather than falling. Gravity seems to disappear inside. Newtonian mechanics explains that elevator, person, and apple all fall with the same acceleration, requiring no extra force and therefore no pressure between them. Gravity remains, pulling everyone alike. Einstein noticed a different reading: within a sufficiently small region, no mechanical experiment distinguishes a freely falling elevator from deep space far from gravity. If they are indistinguishable, perhaps gravity is not an ordinary force but spacetime's geometry. This equivalence principle opens general relativity in Chapter 41. Mass as resistance to acceleration reveals another face here: why is inertial mass, resisting acceleration, always equal to gravitational mass, attracted by Earth? Newton treats their equality as coincidence; Einstein treats them as the same thing. A question hidden in an elevator ultimately overturned gravitational theory.
\end{thought}

\begin{puzzles}
\item A bus traveling forward brakes sharply. A helium balloon on a string floats near the ceiling. At braking, does it drift toward the front or the rear? \emph{Hint}: consider the air first, not the balloon. Inertia carries air forward, making pressure slightly higher at the front and lower at the rear. The balloon behaves as though it dislikes high pressure and prefers low, always squeezed toward lower pressure. Decide where air piles up and the answer follows, exactly opposite your body's lean.
\item A rocket flies in space surrounded by vacuum, with no ground to push against and no air to shove. What does its exhaust flame push on? How can it accelerate? \emph{Hint}: enlarge the system boundary to include rocket and all already expelled gas. The rocket throws gas backward and the gas pushes it forward, with the third law doing its work internally. Consider what happens when you throw carried sandbags hard backward from a boat, one by one.
\item Standing on a scale, you suddenly squat quickly. At the first instant, is the reading greater than, less than, or equal to your true weight? \emph{Hint}: your center of mass initially \emph{accelerates} downward. By \(F_{\text{net}}=ma\), downward acceleration requires downward net force, so gravity exceeds support. The scale displays the support force. Think through direction before answering, then predict its reading as your downward squat comes to a stop.
\item Raw eggs dropped from the same height onto concrete and thick foam usually survive only the foam landing. The velocity change before stopping is nearly the same, so where is the difference? \emph{Hint}: use the seat-belt account: the same velocity change in less time means more force. Whatever factor the foam increases stopping time by, it reduces shell force by. This is also how foam packaging works.
\end{puzzles}

The three rules now stand: motion is unchanged by default, the first law; net force determines momentum change, the second; and forces always occur in pairs on different objects, the third. We have not yet inventoried the force family itself: gravity, normal force, tension, friction, and air resistance. Where do they come from, what rules do they obey, and when do they adjust as needed? That is next chapter's task. Meanwhile, this chapter's repeated phrase, net force changes momentum, will grow in Chapter 10 into an entire accounting system for collisions.
